% Variante : theme=ocots + listing=card (voir la note affichée dans le PDF).
\documentclass[11pt]{ocots-book}
\usepackage[lang=fr, theme=ocots, listing=card, solutions=inline, math=analysis]{ocots}
\lstset{language=[LaTeX]TeX}

\title{Variante d'option}
\author{Prénom \textsc{Nom}}

\newcommand{\ocotsvariantnote}[1]{\begin{quote}\itshape Variante : #1\end{quote}}

\begin{document}
\maketitle

\begin{lstlisting}
\usepackage[..., theme=ocots, listing=card]{ocots}
\end{lstlisting}
\ocotsvariantnote{\texttt{listing=card} (r\'eglage par d\'efaut du th\`eme \texttt{ocots} : la variante force explicitement la valeur).}
% Corps partagé par les variantes de facette (examples/themes/<theme>/variants/) :
% aucune option, aucun préambule. Repris de l'ancien examples/variants/body.tex,
% avec des paires code/rendu et une ligne de polices pour la variante « fonts ».

\lstset{language=[LaTeX]TeX}

\chapter{Chapitre de démonstration}
\begin{chapterintro}
    Une introduction de chapitre, avec \ie un exemple de macro localisée et
    \enquote{un texte cité}. Quelques symboles à police variable :
    $\N$, $\Z$, $\R$, $\calset{O}$.
\end{chapterintro}
\section{Une section}

\begin{lstlisting}
\begin{theorem}{Un theoreme}{demo}
    Enonce.
\end{theorem}
\end{lstlisting}
\begin{theorem}{Un théorème}{demo}
    Énoncé.
\end{theorem}

\begin{lstlisting}
\begin{proof}
    Demonstration.
\end{proof}
\end{lstlisting}
\begin{proof}
    Démonstration.
\end{proof}

\begin{lstlisting}
\begin{definition}{Une definition}{}
    Enonce.
\end{definition}
\end{lstlisting}
\begin{definition}{Une définition}{}
    Énoncé.
\end{definition}

\begin{lstlisting}
\begin{remark}
    Une remarque.
\end{remark}
\end{lstlisting}
\begin{remark}
    Une remarque.
\end{remark}

\begin{lstlisting}
\begin{assumption}
    Une hypothese, etiquetee automatiquement.
\end{assumption}
\end{lstlisting}
\begin{assumption}
    Une hypothèse, étiquetée automatiquement.
\end{assumption}

\begin{lstlisting}
\begin{exercise}[points=4]
    Enonce de l'exercice.
\solution
    Corrige.
\end{exercise}
\end{lstlisting}
\begin{exercise}[points=4]
    Énoncé de l'exercice.
\solution
    Corrigé.
\end{exercise}

\end{document}
